%%%%%%%%%%  scan idx 627  |  folio 614  %%%%%%%%%%
en procédant par récurrence décroissante, supposons qu'elle est connue pour \(n+1\). Quitte à remplacer \(S\) par \(S_{n+1}=S'_{n+1}\), on peut supposer que \(R^if_*(F)\) est \(S\)-plat pour \(i\ge n+1\). Comme \(F\) est plat sur \(S\), il résulte alors de EGA III 6.9.8 que \(R^nf_*(F)\) commute à tout changement de base. Par conséquent la représentabilité de \(S_n\) résulte du lemme suivant (voir : J.P. Murre, Representation of unramified functors, Séminaire Bourbaki~; Mai 1965, p.294-11, th.2) :

\textbf{Lemme 4.5.} \uline{Soient} \(X\to S\) \uline{un morphisme propre de présentation finie,} \(F\) \uline{un faisceau de présentation finie sur} \(X\)\uline{. Alors le sous-foncteur} \(S_F\) \uline{de} \(S\) \uline{qui rend} \(F\) \uline{plat sur} \(S\) \uline{est représentable par un monomorphisme surjectif de présentation finie (qui n'est bien sûr pas une immersion en général).}

Pour représenter la condition 3) de 4.3, on note que d'après 1) et 2), \(f_*(L)\) et \(g_*(p_1^*L)\) commutent aux changements de base \(I\to S\), et donc \(\operatorname{coker}(\pi_1-\pi_2)\) y commute aussi~; ensuite on applique 4.5.

Pour représenter 4), on note que d'après 3) et le fait que \(f_*(L)\) et \(g_*(p_1^*L)\) sont plats et commutent aux changements de base, \(N=\ker(\pi_1-\pi_2)\) commute aussi~; ensuite on applique le

\textbf{Lemme 4.6.} \uline{Soient} \(\lambda:X\to S\) \uline{un morphisme propre de présentation finie,} \(u:F\to G\) \uline{un homomorphisme de} \({O}_X\)\uline{-Modules de présentation finie, avec} \(G\) \uline{plat sur} \(S\)\uline{. Alors le sous-foncteur} \(S_1\) \uline{de} \(S\) \uline{qui rend} \(u\) \uline{un isomorphisme est représentable par un sous-schéma ouvert de} \(S\)\uline{, de présentation finie.}

On se ramène comme d'habitude au cas \(S\) noethérien, en utilisant EGA IV 8 et 11.2.6, ce qui nous dispense de nous préoccuper de la présentation finie. Représentons d'abord le foncteur qui rend \(u\) un épimorphisme,
